\section{Asymptotic midpoint growth}\label{sec:consequences}

The main tail estimate gives a uniform increase in the endpoint norms
on a finite-codimensional subspace. We specify the two standard
midpoint moduli before deriving this consequence for both models.

For an infinite-dimensional real Banach space $Y$ and $t>0$, define
the maximum-based asymptotic midpoint modulus by
\begin{equation}\label{eq:max-modulus}
 \overline\theta_Y(t)=\inf_{\|x\|=1}
 \sup_{F\in\operatorname{cof}(Y)}
 \inf_{y\in F,\,\|y\|=1}
 \left(\max\{\|x+ty\|,\|x-ty\|\}-1\right),
\end{equation}
where $\operatorname{cof}(Y)$ denotes the closed finite-codimensional
subspaces. This is the convention of \cite[Definition~2.2]{DKRRZ2016}.
A norm is AMUC when this modulus is positive for every $t>0$.

We shall also use the averaged convention
\begin{equation}\label{eq:averaged-modulus}
 \widehat\delta_Y(t)=\inf_{\|x\|=1}\sup_{F\in\operatorname{cof}(Y)}
 \inf_{y\in F,\,\|y\|=1}
 \left(\frac{\|x+ty\|+\|x-ty\|}{2}-1\right).
\end{equation}
For every infinite-dimensional real Banach space,
\begin{equation}\label{eq:modulus-comparison}
 \tfrac12\overline\theta_Y(t)\le\widehat\delta_Y(t)
 \le\overline\theta_Y(t).
\end{equation}
The second inequality is pointwise. For the first, fix a unit center
$x$ and a norming functional $f$. Intersect a finite-codimensional
subspace used for the maximum with $\ker f$. Both signed norms are then
at least one, and an excess $a$ for their maximum gives an excess at
least $a/2$ for their average. Arbitrary approximations to the supremum
over subspaces, followed by the infimum over centers, prove the first
inequality. The two conventions therefore agree as qualitative
properties, but their numerical bounds must be distinguished.

For either symmetric expression, allowing $\|y\|\ge1$ in place of
$\|y\|=1$ does not change the inner infimum. Indeed, the maximum and
the average are even convex functions of the scalar multiplying a fixed
direction, and are nondecreasing on the nonnegative half-line. This also
shows that both moduli are nondecreasing in $t$.

\subsection{Averaged midpoint growth at every radius}

Write $\widehat\delta(t)$ for the averaged modulus of the chosen
space, defined in \eqref{eq:averaged-modulus}.

\begin{corollary}\label{mol:cor:all-radius-averaged}
In either of the segment-atomic completions considered above,
\begin{equation}\label{mol:eq:all-radius-averaged}
 \widehat\delta(t)\ge
 \frac{\sqrt{1+t^2/4}-1}{2}\qquad(t>0).
\end{equation}
\end{corollary}

\begin{proof}
For a unit $x=P_Hx$ and $y\in\ker P_H$ with $\|y\|\ge1$,
Theorem~\ref{mol:prop:joint-budget} applied to $ty$ gives
\[
 \max\{\|x+ty\|,\|x-ty\|\}\ge\sqrt{1+t^2/4}.
\]
Both endpoint norms are at least one, by applying the contractive
projection $P_H$.  Their average therefore exceeds one by at least
the right side of \eqref{mol:eq:all-radius-averaged}.

For an arbitrary unit center $x$, choose a finitely supported unit
vector $x_0$ with $\|x-x_0\|<\varepsilon$, and let $H$ contain
its ancestral support.  The estimate for $x_0$ applies to every
$y\in\ker P_H$ of norm at least one.  Replacing $x_0$ by $x$
loses at most $\varepsilon$ in the average.  Take the infimum over
these $y$, then the supremum over finite-codimensional subspaces,
and finally let $\varepsilon\downarrow0$.  The resulting bound
is uniform over the original unit center.
\end{proof}

